\manualchapter{Troubleshooting and further reading}{sec:troubleshooting}
\begin{description}
\item[No picture or no game response.]
Load a supported ROM and release HANG and its gate. Check for a load-error
dialog. A title screen may need Start or Select before the game responds.
\item[The screen moves but there is no sound.]
Reach a part of the game that makes sound. Check the voice levels and Mix
Volume. Voices with patched individual outputs are removed from MIX.
If all five are patched, MIX is silent. Check your downstream audio routing.
\item[A gate does nothing or works only once.]
Reach about 2 V to activate, then return below about 0.1 V to re-arm.
Lengthen very short pulses. Controller actions also depend on game polling;
a held gate is a held button.
\item[Clock CV has no effect.]
The attenuverter starts at zero. Turn it up and keep the base knob away
from its limits. MOD takes a speed voltage, not a synchronization clock.
\item[LOAD does nothing or the loop is inconsistent.]
Make a snapshot first. Loading a new ROM or initializing the module clears
it. Gate patterns and game logic can change what happens after LOAD, and
snapshots do not contain the audio queue or frame-clock phase.
\item[HANG leaves a voltage on the audio or clock jack.]
This is the implemented behavior. Use an external VCA or mute if you need
silence while the machine is held.
\item[The saved patch cannot find its game.]
Restore the same ROM to the stored path and reopen the patch. Loading a
ROM manually starts a fresh session and clears the snapshot slot.
\item[High-speed operation is expensive or sounds rough.]
Reduce Clock Speed and its modulation depth. The implementation does more
CPU emulation work per Rack sample at high rates; its audio conversion does
not guarantee smooth resampling across speed changes.
\end{description}

\subsection{Sources and acknowledgments}
This manual follows the module controls, emulator, audio wrapper, and theme
handling in the repository's \texttt{src/} directory. The maintainer's
\href{https://github.com/Kautenja/RackNES/tree/master/manual}{manual guide}
maps each topic to its implementation.
RackNES builds on Amish Naidu and contributors' SimpleNES; Shay Green's
(blargg's) Nes\_Snd\_Emu, Blip\_Buffer, and nes\_ntsc; and Ren\'e
Nyffenegger's cpp-base64. The
\href{https://github.com/Kautenja/RackNES/blob/master/LICENSING.md}{license guide}
links the bundled component notices and license texts. CV Genie was developed
by \textbf{@anlexmatos}. The companion
\href{https://github.com/Kautenja/RackNES/releases/latest/download/CVGenie.pdf}{CV Genie manual}
covers its eight input rows and memory maps.
